\subsection{Reference protection and independent charging input}
\label{sec:reference-protection}

The matrix error comparisons motivate protection relative to a common
executable route.  We fixed the rule in \eqref{eq:protected-route} before
running its validation and evaluation.  Validation used 512 paired 2022
scenarios in each of seven conditions and compared
$\nu\in\{0,0.00025,0.0005,0.001\}$.  A candidate margin first had to limit
its original-condition mean cost increase over central PC to 0.05\% of
matched cost.  Among passing margins, selection minimized mean positive
excess above matched across the six perturbed conditions, with ties assigned
to the smaller margin.  Both route portfolios selected $\nu=0.001$.
Table~\ref{tab:guarded-validation} retains all validation choices.
The subsequent 2023 evaluation used 4,096 original-condition pairs and
1,024 pairs in each of the six perturbed conditions and electrical stress.
Because this year had already been examined during revision, it was a
fixed-rule reevaluation rather than an unused test year.

Let $E^+(P)=n^{-1}\sum_i[Q_i(P)-Q_i(P_b)]_+$, where each $P_b$ is the
matched route for that scenario and portfolio.  This quantity measures
additional cost incurred when a method performs worse than its reference.
It differs from the upper tail of total cost, which also includes scenario
severity shared by every policy.  For the aggregate error comparison, we
first averaged the six condition-specific positive-excess differences
within each of the 1,024 scenario identifiers.  Resampling those identifiers
preserved dependence between the conditions.  Mean intervals used 20,000
paired resamples.  The two portfolio comparisons received Holm adjustment,
while all 80 additional mean-cost comparisons formed a separate family.
The original-condition noninferiority check required a one-sided 97.5\%
bootstrap upper bound below zero for
$Q_i(P_g)-Q_i(P_c)-0.0005Q_i(P_b)$.

\begin{table}[t]
\centering
\caption{Reference protection on the original common routes. $G$ denotes PC with reference protection and $C$ central PC. $E^+$ is the mean positive excess cost above forecast matched. Intervals are paired marginal 95\% intervals. The original condition has 4,096 pairs and each other condition has 1,024 pairs.}
\label{tab:guarded-primary}
\footnotesize\setlength{\tabcolsep}{3pt}\renewcommand{\arraystretch}{1.12}
\begin{tabular}{@{}p{.22\columnwidth}p{.40\columnwidth}rr@{}}
\toprule
Condition & $G-C$ [95\% CI] & $E_C^+$ & $E_G^+$ \\
\midrule
Original & 0.0375 [-0.2279, 0.3017] & 1.0834 & 0.0657 \\
$0.75M$ & -3.0865 [-4.1245, -2.1035] & 5.5237 & 2.0125 \\
$1.25M$ & 0.6297 [0.0687, 1.2777] & 0.6728 & 0.0689 \\
Noise 0.15 & 0.9903 [0.5097, 1.5282] & 0.2862 & 0.0450 \\
Couplings doubled & -1.2044 [-2.0505, -0.3899] & 2.5848 & 0.1524 \\
Couplings zero & 0.4804 [0.1196, 0.8722] & 0.4702 & 0.0197 \\
Nonlinear & 0.7658 [0.4804, 1.0707] & 0.2579 & 0.0088 \\
Electrical stress & -0.6917 [-5.9540, 4.6563] & 10.7569 & 1.3159 \\
\bottomrule
\end{tabular}
\end{table}


On the original routes, reference protection retained a mean cost of
1112.5728, compared with 1112.5353 for central PC and 1113.6789 for matched.
The protected minus matched difference was $-1.1061$, with interval
$[-1.3760,-0.8582]$ and Holm adjusted paired $t$ probability
$3.93\times10^{-15}$.  The reduction was 0.0993\%.  Its mean penalty relative
to central PC was 0.0375, or 0.0034\%, and the noninferiority upper bound
was $-0.2550$.  The rule accepted a replacement in 6.20\% of scenarios.
Consequently, most routes remained matched, but the replacements retained
most of central PC's average benefit.  The fraction of scenarios worse than
matched fell from 11.94\% for central PC to 0.90\%, and $E^+$ fell from
1.0834 to 0.0657.  These frequencies describe realized outcomes rather than
an unconditional guarantee from the model score.

Across the six error conditions, the paired decrease in positive excess was
1.2481, with interval $[-1.4785,-1.0296]$ for protected minus central and
Holm adjusted probability $6.80\times10^{-26}$.
Table~\ref{tab:guarded-primary} shows the individual conditions as well.
Protection reduced additional losses under an underscaled matrix and
expanded unobserved couplings, but it gave up some of central PC's average
gain under an enlarged matrix, coefficient noise, zero couplings, and
nonlinear propagation.  Positive-excess protection is therefore a selective
tradeoff, not uniform mean superiority.  Nor does it imply lower total-cost
tails.  The original-condition upper 5\% mean cost exceeded central PC by
2.7673, with exploratory interval $[0.3615,5.5575]$.  Both portfolio tail
comparisons are retained in Table~\ref{tab:guarded-tails}.

The separate parameter-independent route construction tests whether a
favorable score ablation could arise from using unknown parameters to
generate the routes.  Table~\ref{tab:guarded-ablation} reports all conditions
on this common alternative portfolio.  Removing the five unsupported terms
from the score increased original-condition cost over central PC by 0.9937,
with interval $[0.7761,1.2289]$ and Holm adjusted probability
$1.26\times10^{-15}$.  Thus simply deleting those terms did not resolve
parameter uncertainty.  Reference protection instead increased mean cost by
0.1039 while passing the fixed noninferiority check, whose upper bound was
$-0.3419$.  Its aggregate positive-excess decrease was 0.1823, with interval
$[-0.2691,-0.1089]$ and Holm adjusted probability $1.01\times10^{-5}$.
Its mean benefit over matched did not survive the full comparison correction.
These results support the protection mechanism on a second route construction
without attributing differences between the portfolios to a single score.

An independent geographic input used the City of Palo Alto charging
records.\footnote{City of Palo Alto, Electric Vehicle Charging Station Usage,
published February 2021. The source record is available at
\url{https://data.paloalto.gov/dataviews/257812/electric-vehicle-charging-station-usage-july-2011-dec-2020/}.}
The complete municipal source contained 259,415 sessions.  Fixed quality
rules and a requirement of at least 100 valid sessions in 2017 and 2018
retained 33 stations and 159,587 sessions during 2017 to 2020, including
20,042 in the 2020 evaluation year.  Arrival counts and session energy were
observed.  Hourly energy was reconstructed uniformly over the reported
charging duration, giving 1,416,780.194 kWh and conserving retained energy
to numerical precision.  Station eligibility, coordinates, mean demand,
and budgets used 2017 and 2018 only.  The charging, communication, and
repair priority budgets were 54.4678 kWh per hour, 50.3244 service units,
and 10 priority units, with the same construction rules as Boulder.

The restoration parameters and selected margin were unchanged.  One demand
input used the selected Boulder graph model with identical learned weights
and normalization, replacing only the nonlearned adjacency by the Palo Alto
station graph.  This was a retrospective geographic transfer, not a claim
that the later-trained model had been deployed in 2020.  The second input
used the corresponding hours of the preceding week.  Both forecasts were
fixed before computing external policy outcomes and were evaluated separately.
The electrical feeder, abstract transition, communication conditions, and
repair duration model were retained, while route distances used the new
station coordinates.  Constructed packet response functions were assigned
by station index.  The new observations therefore test charging demand and
station geometry rather than a jointly observed local infrastructure event.

\begin{table}[t]
\centering
\caption{Independent municipal charging input with 1,024 pairs per row. Original and independent identify the two common route portfolios. Reduction is relative to forecast matched. The final column is the one-sided 97.5\% upper bound after subtracting the fixed 0.05\% noninferiority margin.}
\label{tab:guarded-transfer}
\footnotesize\setlength{\tabcolsep}{3pt}\renewcommand{\arraystretch}{1.12}
\begin{tabular}{@{}p{.24\columnwidth}p{.38\columnwidth}rr@{}}
\toprule
Forecast and routes & $G-C$ [95\% CI] & Red. (\%) & NI upper \\
\midrule
Graph, original & -0.2826 [-0.6408, 0.0433] & 0.1806 & -0.2173 \\
Graph, independent & -0.0227 [-0.1109, 0.0607] & 0.0124 & -0.1828 \\
Weekly, original & -0.1761 [-0.4869, 0.1121] & 0.1338 & -0.1487 \\
Weekly, independent & -0.0735 [-0.1912, 0.0278] & 0.0169 & -0.2193 \\
\bottomrule
\end{tabular}
\end{table}


Each input used 1,024 paired scenarios per condition, with all eight
conditions and both portfolios retained.  External mean comparisons formed
one 160-test Holm family.  On the original routes, reference protection
reduced matched cost by 0.1806\% with the graph forecast and 0.1338\% with
weekly persistence, with adjusted probabilities 0.00481 and 0.000264.
The corresponding differences from central PC were $-0.2826$ and $-0.1761$,
but their intervals included zero.  Table~\ref{tab:guarded-transfer} shows
that both fixed inputs and both route portfolios passed the predeclared
noninferiority check. Tables~\ref{tab:guarded-external-published} and
\ref{tab:guarded-external-independent} give the individual error and stress
conditions under both forecasts. The validation, mean-cost, error-loss, and external
conditions thus supported adopting reference protection, while retaining
central PC and worst absolute-score PC as distinct comparisons.
